\documentclass[a4paper,12pt]{scrartcl}
\usepackage[singlespacing]{setspace}

% Layout
\usepackage[a4paper, left=2.5cm, right=2cm, top=2cm, bottom=3.2cm]{geometry} % kleinere Ränder
\usepackage{parskip}

% Umlaute in der Datei erlauben, auf Deutsch umstellen
\usepackage{fontspec}
\usepackage{lmodern}
\usepackage[english]{babel}

% Mathesymbole und Ähnliches
\usepackage{amsmath}
\usepackage{amsthm}
\usepackage{mathtools}
\usepackage{amssymb}
\usepackage{microtype}
\usepackage{stmaryrd}
\usepackage{MnSymbol}
\usepackage{bussproofs}		% Sequenzenkalkül
\usepackage{xcolor}

% Grafiken und PDFs einfügen
\usepackage{graphicx}
\usepackage{pdfpages}
\usepackage{float}

% PDF-Tools
\usepackage[hidelinks, unicode]{hyperref}

% Abbildungen
\usepackage{tikz}
\usetikzlibrary{arrows,calc,positioning}
\tikzset{ 
   state/.style={circle,fill,inner sep=2pt,label=below:#1\strut},
   arrow/.style={->, shorten <=3pt, shorten >=3pt, >=stealth',thick},
   knoten/.style={draw, circle, minimum size=8mm}} 
   
% Infoboxen 
\definecolor{orange}{RGB}{230,159,0}
\definecolor{darkgray}{RGB}{70,90,70}

\usepackage{tcolorbox}

\newenvironment{LP}[1]{%
\begin{minipage}{0.48\textwidth}
\begin{tcolorbox}[colback=orange!20!white, colframe=orange!80!black, title={\sffamily \textbf{#1}}]
}{
\end{tcolorbox}
\end{minipage}}

\newenvironment{alg}[1]{%
\begin{tcolorbox}[colback=red!15!white, colframe=red!85!black, title={\sffamily \textbf{#1}}]
}{
\end{tcolorbox}}

\newenvironment{info}[1]{%
\begin{tcolorbox}[colback=darkgray!15!white, colframe=darkgray!75!black, title={\sffamily \textbf{#1}}]
}{
\end{tcolorbox}}

\newcommand{\ein}{\textbf{Eingabe: }}
\newcommand{\reg}{\textbf{Bedingungen:}}
\newcommand{\aus}{\textbf{Ziel: }}

% Reelle, Natürliche, Ganze, Rationale Zahlen
\newcommand{\R}{\ensuremath{\mathbb{R}}}
\newcommand{\N}{\ensuremath{\mathbb{N}}}
\newcommand{\Z}{\ensuremath{\mathbb{Z}}}
\newcommand{\Q}{\ensuremath{\mathbb{Q}}}

% Fraktur für Strukturen
\newcommand{\A}{\ensuremath{\mathfrak A}}
\newcommand{\B}{\ensuremath{\mathfrak B}}
\newcommand{\I}{\ensuremath{\mathfrak I}}

% Körper
\newcommand{\K}{\ensuremath{\mathcal{K}}}
%\newcommand{\G}{\ensuremath{\mathcal{G}}}

% Makros für logische Operatoren
\newcommand{\xor}{\ensuremath{\oplus}} % exklusives oder
\newcommand{\impl}{\ensuremath{\rightarrow}} % logische Implikation
\newcommand{\Lor}{\ensuremath{\bigvee}}
\newcommand{\Land}{\ensuremath{\bigwedge}}

% Meistens ist \varphi schöner als \phi, genauso bei \theta
\renewcommand{\phi}{\varphi}
\renewcommand{\theta}{\vartheta}

\newcommand*{\todo}{\textbf{\textcolor{red}{TODO}}}		% big, red, apparent TODO

% SEQUENZENKALKÜL
\usepackage{bussproofs}
\newcommand{\yields}{\ensuremath{\mathrel{\Rightarrow}}}

% Aufzählungen anpassen (alternativ: \arabic, \alph)
\renewcommand{\labelenumi}{(\roman{enumi})}

% BITTE NICHT ÄNDERN: interne Kommandos für die Informationen (Blattnummer, Gruppe, ...)
\newcounter{sheet}
\newcommand{\printsheet}{\arabic{sheet}}

\newcommand{\printgroup}{?}
\newcommand{\group}[1]{%
\renewcommand{\printgroup}{#1}%
}

\newcommand{\printmembers}{}
\newcommand{\printmember}[3]{{#2} {#3} & {#1} \\}

% Für Vektoren
\newcommand{\myvec}[1]{\ensuremath{\begin{pmatrix}#1\end{pmatrix}}}
\def\spvec#1{\left(\vcenter{\halign{\hfil$##$\hfil\cr \spvecA#1;;}}\right)}
\def\spvecA#1;{\if;#1;\else #1\cr \expandafter \spvecA \fi}

\newtheorem{proposition}{Proposition}[sheet]

\newcommand{\pdfmembers}{}
\newcommand{\pdfmember}[3]{{#1} {#3}, {#2}; }

\newcommand{\member}[3]{%
\expandafter\renewcommand\expandafter\printmembers\expandafter{\printmembers\printmember{#1}{#2}{#3}}%
\expandafter\renewcommand\expandafter\pdfmembers\expandafter{\pdfmembers\pdfmember{#1}{#2}{#3}}%
}

\AtBeginDocument{\hypersetup{
    pdftitle = {Übungsblatt \printsheet},
    pdfauthor = {Abgabegruppe \printgroup: \pdfmembers},
    pdfsubject = {Online Algorithms}
}}

% <-------- HIER alle Informationen eintragen ========================
\setcounter{sheet}{1} % Nummer des Blatts

% alle Gruppenmitglieder in der Form \member{Matrikelnummer}{Vorname}{Nachname}
\member{5481114}{Danje}{Petersen}

\begin{document}

% Kopfzeile
{\raggedright
\begin{tabular}{l}
    Online Algorithms \\
    SS 2024 \\
\end{tabular}}
\hfill
{\Large Sheet 0\printsheet}
\hfill
\begin{tabular}{l l}
\today{} \\
\printmembers 
\end{tabular}
\hrule


% <-------- HIER beginnt die Lösung ========================

\section*{Exercise \printsheet.1}
We construct a directed graph $G=(V,E,c)$ that models the problem. For each travel request there is one vertex $v_i$ and there is one starting and final vertex $s,t$. Therefore:
\[V=\{s, v_1, \dots v_n, t\}\]
For each travel request $(t_i,c_i)$ there exist two edges starting in its vertex $v_i$. One models the decision not to buy a BahnCard and connects $v_i$ with $v_{i+1}$ and has costs $c_i$. Additionally, there is an edge from that request to the request that has the lowest time $t_j$ with $t_j > t_i+T$. This models the decision to buy a BahnCard at $t_i$. The edge has the costs of each of the travel requests that will be discounted multiplied by $\beta$.
\begin{align*}E= & \{(v_i,v_{i+1}) \mid i \in \{1,\dots ,n\} \text{ with } c(v_i,v_{i+1}) = c_i \} \\
&\cup \{(v_i,v_j) \mid i \in \{1,\dots ,n\} \text{ where } t_j \text{ is next request after } t_i+T \\
& \text{ with } c(v_i,v_j) = \sum_{k:t_k \in [t_i,t_i+T)} \beta \cdot c_k \}\end{align*}

On this directed graph we do a search for the shortest $s-t-$path which is possible in $O(n)$ since there are at most $2n$ unique edges in the graph. Calculating the edge costs is possible in $O(n)$ when going through the travel requests chronologically. \\
The shortest path now corresponds to the optimal decisions and its costs to the costs of the bought tickets.


\section*{Exercise \printsheet.2}
\subsection*{(a)}
\begin{proposition} For each phase $[\tau_i,\tau_{i+1})$ with $1 \leq i\leq k$ the interval $[\tau_i,\tau_i+T)$ is expensive. \end{proposition}
\begin{proof}
At $\tau_i$ OPT buys a BahnCard. Assume, $[\tau_i,\tau_i+T)$ is not expensive. Then it is cheaper not to buy a BahnCard than to buy one. Therefore OPT would not buy a BahnCard at $\tau_i$. This contradicts the fact that OPT buys a BahnCard at $\tau_i$.
\end{proof}

\begin{proposition} Any subinterval of $[\tau_i+T,\tau_{i+1})$ of length smaller than $T$ is cheap. \end{proposition}
\begin{proof}
Assume $[\tau_i+T,\tau_{i+1})$ is expensive. Then it would be optimal to buy a BahnCard at the time of the next travel request $t_j$ after $\tau_i+T$. Since OPT makes the optimal choice, it would buy a BahnCard at $t_j$. Since $t_j \notin \{\tau_1, \dots \tau_k\}$, OPTdoes not buy a BahnCard. Therefore $[\tau_i+T,\tau_{i+1})$ can't be expensive.  
\end{proof}

\pagebreak

\subsection*{(b)}
\begin{proposition} For $I=[0,\tau_1)$, it holds that $c_{SUM} \leq c_{OPT}$. \end{proposition}
\begin{proof}
Since OPT buys its first BahnCard at $\tau_1$, $c_{OPT} = \sum_{i: t_i \in I}c_i$ and $I$ can't be expensive, thus $\sum_{i: t_i \in I}c_i < c^*$. SUM buys a BahnCard as soon as for $j$ the sum $\sum_{k:t_k \in [t_j-T,t_j)} c_k \geq c^*$. Since this can't be true at any time in $I$, SUM does not buy a BahnCard in $I$. Therefore
\[c_{SUM} = \sum_{i: t_i \in I}c_i = c_{OPT} \Rightarrow c_{SUM} \leq c_{OPT} \text{ for } I = [0,\tau_1).\]
\end{proof}


\subsection*{(c)}
If \textit{SUM} doesn't buy a BahnCard, then $c_{SUM} = \sum_{j:t_j \in I}c_j$.
For \textit{OPT} the costs would be 
\[c_{OPT} = C + \sum_{j:t_j \in [\tau_i,\tau_i+T)} \beta c_j + \sum_{j:j_j \in [\tau_i+T,\tau_{i+1})}c_j\]
The worst case for \textit{SUM} would be that \textit{OPT} only pays discounted prices and therefore
\[c_{OPT} = C + \sum_{j:t_j \in [\tau_i,\tau_i+T)} \beta c_j = C + \beta c_{SUM}\]
This scenario should not be possible, since whenever \textit{OPT} buys a BahnCard, there needs to be an expensive interval in $[\tau_i,\tau_{i+1})$ at which's end, \textit{SUM} would buy a BahnCard.


\subsection*{(d)}
We partition the interval $I$ into three intervals $I_1,I_2,I_3$. $I_1$ is the timeframe in which \textit{SUM} still owns a BahnCard that was bought in the time before $\tau_i$. $I_2$ is the interval in which \textit{SUM} doesn't own a BahnCard and $I_3$ the interval starting when \textit{SUM} buys a new BahnCard.
For each interval we denote the costs of the travel requests in it as:
\[C_j = \sum_{k: t_k \in I_j} c_k\] 
Since \textit{OPT} buys a BahnCard at $\tau_i$ that expires at $\tau_i + T$, the interval $[\tau_i,\tau_i+T)$ must be expensive. Since \textit{SUM} buys a BahnCard whenever the current request makes the previous interval expensive, it buys a BahnCard at $t'$ which is the time of the last travel request $(t',c')$ before $\tau_i+T$. Therefore $I_3 = [t',\tau_i+1)$\\
We can assume w.l.o.g. that there are no non-discounted requests for \textit{SUM} in $I_3$, since both \textit{SUM} and \textit{OPT} would have to pay the full price. 
The worst case for \textit{SUM} would be, if $C_1$ and $C_3$ would be very low and $C_2$ very high, since $I_2$ is the only interval in which only \textit{OPT} gets the discount.

\begin{align*}
\frac{c_{\textit{SUM}}}{c_{\textit{OPT}}} &= \frac{C + \beta C_1 + C_2 + \beta C_3}{C + \beta C_1 + \beta C_2 + \beta c' + C_3 - c'} \\
&= \frac{C + \beta C_1 + C_2 + \beta C_3}{\frac{C}{1-\beta} + C_3 - c'} \\
&\leq  \frac{C + C_1 + C_2 + C_3}{\frac{C}{1-\beta} + C_3 - c'} \quad \mid \text{in the worst case, \textit{SUM} has to pay almost all the full prices}\\
&= \frac{C+C_2}{\frac{C}{1-\beta}} \quad \mid \text{Let }C_1, C_3 \text{ be very small}\\
&= \frac{(1-\beta)(C+C_2)}{C} \\
&= \frac{(1-\beta)C}{C} + \frac{(1-\beta)C_2}{C} \\
&\leq  1-\beta + 1 \quad \mid C \text{ must be smaller than the saved costs} \\
&= (2-\beta)
\end{align*}


\section*{Exercise \printsheet.3}
\subsection*{(a)}
\begin{proposition} A is $c$-competitive, i.e. $\sum_{i=1}^n c_A(i) \leq c \cdot \sum_{i=1}^n c_{OPT}(i)$, if for every request sequence $r$, there is a potential function $\Phi_r$ such that $c_A(i) + \Phi_r(i) - \Phi_r(i-1) \leq c \cdot c_{OPT}(i)$. \end{proposition}
\begin{proof}
\begin{align*}
\sum_{i=1}^n c_A(i) & \leq \sum_{i=1}^n c_A(i) + \sum_{i=1}^n \Phi_r(i) - \Phi_r(i-1) \mid \text{ since } \Phi_r \text{ can't be negative} \\
& = \sum_{i=1}^n c_A(i) + \Phi_r(i) - \Phi_r(i-1) \\
& \leq \sum_{i=1}^n c \cdot c_{OPT}(i) \\
& = c \cdot \sum_{i=1}^n c_{OPT} (i)
\end{align*}
\end{proof}

\pagebreak

\subsection*{(b)}
The algorithm doubles the buffer space each time the limit is reached. \\
For each element $s_i$ of $s = s_1,\dots s_n$ the algorithm determines if the buffer is full. If it is full, a new buffer is allocated with double the size of the old buffer. If it is not full, $s_i$ is written into the buffer. \\

In the worst case, the last element of $s$ is added when the buffer is full and needs to be reallocated. In this case the size of the buffer is $2(n-1) \leq 2 |s|$.

To prove that the algorithm is $3$-competitive, we use a potential function $\Phi$.

\[\text{For every }s_i \in s: \Phi(i) = \begin{cases} \Phi(i-1)+2 & \text{if buffer is not full} \\ 0 & \text{otherwise} \end{cases} \quad \text{ with } \quad \Phi(0)=0\]

Since no element of $\Phi$ is negative, this is a valid potential function.

\begin{proposition} $\forall s_i \in s: c_A(i) + \Phi(i) - \Phi(i-1) \leq 3 \cdot c_{OPT}$ \end{proposition}
\begin{proof}
For every element in $s$, $c_{OPT}$ has costs of $1$ to add it to the buffer.
For any $i$ for which the buffer doesn't need to be reallocated the algorithm needs one step, i.e. $c_A(i) = 1$ and therefore the following holds:
\begin{align*}
c_A(i) + \Phi(i) - \Phi(i-1) &= 1 + \Phi(i-1) - \Phi(i-1) + 2 \\
&= 3 \\
&\leq 3 \cdot 1 = 3 \cdot c_{OPT}(i)
\end{align*}

For any $i$ where the buffer needs to be reallocated the cost is $c_A(i) = i$ since $(i-1)$ elements are reallocated and then the new element is added. Since the potential function is \glqq reset\grqq{} every $2^k+1$ elements when the buffer needs to be reallocated, when it happens again, $\Phi(i-1)$ is equal to $2\cdot \left(\frac{(i-1)}{2}-1\right)$. For example, when the $9th$ element is added, the buffer needs to be reallocated. $\Phi$ has last been reset to $0$ at element $5$. Therefore, for each element between $5$ and $9$ the potential function has increased by $2$ which leads to 
\[\Phi(9-1) = \Phi(8) = 2 \cdot \frac{8}{2} -1 = 2 \cdot 3 = 6.\] 

Now, for each element at which the buffer needs to be reallocated it holds that:
\begin{align*}
c_A(i) + \Phi(i) - \Phi(i-1) &= i + 0 - 2 \cdot \left(\frac{(i-1)}{2} - 1\right) \\
&= i - (i-1) + 2 \\
&= 3 \\
&\leq 3 \cdot 1 = 3 \cdot c_{OPT}(i)
\end{align*}
\end{proof}

Combining Proposition 1.4 and Proposition 1.5, it is clear that the algorithm is $3$-competitive.





\end{document}
